\documentclass[11pt]{article}
\usepackage[letterpaper,margin=0.83in,headheight=16pt]{geometry}
\usepackage[T1]{fontenc}
\usepackage[utf8]{inputenc}
\usepackage{lmodern}
\usepackage{microtype}
\usepackage{amsmath,amssymb}
\usepackage{xcolor}
\usepackage{hyperref}
\usepackage{enumitem}
\usepackage{booktabs,longtable,tabularx,array}
\usepackage{fancyhdr}
\usepackage{titlesec}
\usepackage{listings}
\usepackage{tikz}
\usepackage{float}
\definecolor{suiteblue}{HTML}{17365D}
\definecolor{suitelight}{HTML}{EAF1F8}
\definecolor{suitegray}{HTML}{556070}
\hypersetup{colorlinks=true,linkcolor=suiteblue,urlcolor=suiteblue,breaklinks=true}
\titleformat{\section}{\Large\bfseries\color{suiteblue}}{\thesection}{0.55em}{}
\titleformat{\subsection}{\large\bfseries\color{suiteblue}}{\thesubsection}{0.55em}{}
\setlist[itemize]{itemsep=0.25em,topsep=0.35em,leftmargin=1.4em}
\setlist[enumerate]{itemsep=0.25em,topsep=0.35em,leftmargin=1.6em}
\setlength{\parindent}{0pt}
\setlength{\parskip}{0.47em}
\emergencystretch=2em
\lstset{basicstyle=\ttfamily\footnotesize,breaklines=true,frame=single,columns=fullflexible,showstringspaces=false}
\pagestyle{fancy}\fancyhf{}
\lhead{Game Design \& Engineering}\rhead{Technical Reference}\cfoot{\thepage}
\begin{document}
{\huge\bfseries\color{suiteblue} Formal Game Mechanics Reference}\par
\vspace{0.15in}
{\normalsize\color{suitegray} Original engineering synthesis; Schell-aligned concepts explicitly identified}\par
\vspace{0.15in}\hrule\vspace{0.15in}

\section{Source and engineering boundary}
Schell (third edition, 2020 copyright) supplies the conceptual design framing. Formal notation, software architectures, validation gates, metrics, consent protocols, and recommended tooling are supplementary engineering practices introduced by this suite. These are not quotations or source-prescribed standards.


\section{Transition-system notation}
A deterministic discrete mechanic is modeled as $M=(S,A,T,s_0,G)$ with state space $S$, legal action relation $A(s)$, transition $T:S\times A\rightarrow S$, initial state $s_0$, and goal states $G\subseteq S$. Put input validation before applying transitions.


\subsection{Operational review}
\begin{enumerate}
\item What decision or player-facing problem is this technique meant to resolve?
\item Which inputs, assumptions, and environmental constraints determine validity?
\item What test or observation could falsify the intended outcome?
\item What negative or unintended consequences should reviewers document?
\end{enumerate}

\section{Core gameplay loop}
Input collection, command validation, fixed simulation step, event generation, state update, feedback presentation, and durable observations must be distinguished; only selected systems require deterministic lockstep.


\subsection{Operational review}
\begin{enumerate}
\item What decision or player-facing problem is this technique meant to resolve?
\item Which inputs, assumptions, and environmental constraints determine validity?
\item What test or observation could falsify the intended outcome?
\item What negative or unintended consequences should reviewers document?
\end{enumerate}

\section{Finite state machines}
Represent game modes explicitly; make invalid transitions impossible or reject them with a documented error. Prefer exhaustive tests for small state spaces.


\subsection{Operational review}
\begin{enumerate}
\item What decision or player-facing problem is this technique meant to resolve?
\item Which inputs, assumptions, and environmental constraints determine validity?
\item What test or observation could falsify the intended outcome?
\item What negative or unintended consequences should reviewers document?
\end{enumerate}

\section{Time and reproducibility}
With a fixed step $\Delta t$, update discrete simulation at $t_k=k\Delta t$. Rendering may interpolate between the last two states; never silently equate variable frame rates with simulation accuracy.


\subsection{Operational review}
\begin{enumerate}
\item What decision or player-facing problem is this technique meant to resolve?
\item Which inputs, assumptions, and environmental constraints determine validity?
\item What test or observation could falsify the intended outcome?
\item What negative or unintended consequences should reviewers document?
\end{enumerate}

\section{State invariants}
Examples: health is nonnegative; inventory counts remain integer; resources never increase from a pure spend action; outcomes of seeded draws are reproducible. Distinguish deterministic mechanics from intentionally stochastic rules.


\subsection{Operational review}
\begin{enumerate}
\item What decision or player-facing problem is this technique meant to resolve?
\item Which inputs, assumptions, and environmental constraints determine validity?
\item What test or observation could falsify the intended outcome?
\item What negative or unintended consequences should reviewers document?
\end{enumerate}

\section{Graph-based puzzle rules}
Let nodes represent reachable puzzle configurations and edges legal actions. Breadth-first search proves shortest-path reachability on a finite unweighted graph; it does not prove a human player will understand or enjoy the solution.


\subsection{Operational review}
\begin{enumerate}
\item What decision or player-facing problem is this technique meant to resolve?
\item Which inputs, assumptions, and environmental constraints determine validity?
\item What test or observation could falsify the intended outcome?
\item What negative or unintended consequences should reviewers document?
\end{enumerate}

\section{Input and feedback contracts}
Record action names, admissible states, response deadlines, auditory/visual alternatives, cancellation semantics, and accessibility contingencies.


\subsection{Operational review}
\begin{enumerate}
\item What decision or player-facing problem is this technique meant to resolve?
\item Which inputs, assumptions, and environmental constraints determine validity?
\item What test or observation could falsify the intended outcome?
\item What negative or unintended consequences should reviewers document?
\end{enumerate}

\section{Reference evidence register}
\begin{tabularx}{\linewidth}{@{}p{0.13\linewidth}X X@{}}\toprule ID & Evidence or design decision & Verification method\\\midrule
R-01 & \rule{0.95\linewidth}{0.4pt} & \rule{0.95\linewidth}{0.4pt}\\[0.4cm]
R-02 & \rule{0.95\linewidth}{0.4pt} & \rule{0.95\linewidth}{0.4pt}\\[0.4cm]
R-03 & \rule{0.95\linewidth}{0.4pt} & \rule{0.95\linewidth}{0.4pt}\\\bottomrule\end{tabularx}

\end{document}
